\section*{الفصل \ref{ch:g8:combustion-and-fuels} --- الاحتراق: أنواع الوقود ونواتجه وغازات الدفيئة}

\begin{solution}{exo:g8:combustion-and-fuels:1}
\ce{C + O2 -> CO2}.
\end{solution}

\begin{solution}{exo:g8:combustion-and-fuels:2}
إذا كفى الأكسجين صار كل الكربون ثنائي أكسيد الكربون (\omterm{def:g8:combustion-and-fuels:complete}{احتراق تام}). وإذا
قلّ صار جزء منه أول أكسيد الكربون أو سخامًا (\omterm{def:g8:combustion-and-fuels:complete}{احتراق غير تام}).
\end{solution}

\begin{solution}{exo:g8:combustion-and-fuels:3}
فهو لا لون له ولا رائحة ولا طعم: لا ينتبه إليه أحد، ومع ذلك فهو
سام.
\end{solution}

\begin{solution}{exo:g8:combustion-and-fuels:4}
أحفوري: الفحم، والغاز الطبيعي، والبنزين. متجدد: الخشب، والغاز الحيوي،
والإيثانول المصنوع من الشمندر السكري.
\end{solution}

\begin{solution}{exo:g8:combustion-and-fuels:5}
\ce{C3H8 + 5O2 -> 3CO2 + 4H2O}.
\end{solution}

\begin{solution}{exo:g8:combustion-and-fuels:6}
الكربون: 2 \ce{CO2}؛ الهيدروجين: 3 \ce{H2O}؛ والأكسجين في اليمين $4 + 3 =
7$، منها 1 يأتي من الإيثانول، إذن 3 \ce{O2}:
\ce{C2H6O + 3O2 -> 2CO2 + 3H2O}.
\end{solution}

\begin{solution}{exo:g8:combustion-and-fuels:7}
سخام، أي كربون. \omterm{def:g8:combustion-and-fuels:complete}{فالاحتراق غير تام}: اللهب ينقصه
الأكسجين.
\end{solution}

\begin{solution}{exo:g8:combustion-and-fuels:8}
نحو 339 ppm سنة 1980 ونحو 370 ppm سنة 2000: أي ارتفع بنحو 30 ppm.
\end{solution}

\begin{solution}{exo:g8:combustion-and-fuels:9}
كربون الخشب أخذته الشجرة من \omterm{def:g4:air-a-mixture-of-gases:air}{الهواء} قبل بضع سنوات أو عقود: فحرقه يعيد ثنائي
أكسيد الكربون ذاك. أما كربون الفحم فكان مخزّنًا تحت الأرض ملايين السنين:
فحرقه يضيف إلى \omterm{def:g4:air-a-mixture-of-gases:air}{الهواء} ثنائي أكسيد كربون
جديدًا.
\end{solution}

\begin{solution}{exo:g8:combustion-and-fuels:10}
\ce{2C + O2 -> 2CO}.
\end{solution}

\begin{solution}{exo:g8:combustion-and-fuels:11}
$(427 - 316)/316 = 111/316 \approx 0.35$: أي نحو \qty{35}{\percent}.
\end{solution}

\begin{solution}{exo:g8:combustion-and-fuels:12}
في البداية يكفي الأكسجين: \ce{CH4 + 2O2 -> CO2 + 2H2O}. ومع استهلاك
أكسجين الغرفة المغلقة يحترق اللهب وينقصه \omterm{def:g4:air-a-mixture-of-gases:air}{الهواء}:
\ce{2CH4 + 3O2 -> 2CO + 4H2O}. فيتراكم في الغرفة أول أكسيد الكربون،
وهو عديم الرائحة وسام.
\end{solution}

\begin{solution}{pb:g8:combustion-and-fuels:1}
\textbf{1.} \ce{C3H8 + 5O2 -> 3CO2 + 4H2O}.

\textbf{2.} اليسار: 3 C، و8 H، و10 O. اليمين: 3 C، و$4 \times 2 = 8$ H،
و$3 \times 2 + 4 = 10$ O. موزونة.

\textbf{3.} \ce{2C3H8 + 7O2 -> 6CO + 8H2O}.

\textbf{4.} \omterm{def:g8:combustion-and-fuels:complete}{الاحتراق التام}: 5 \omterm{def:g7:atoms-and-molecules:molecule}{جزيئات} أكسجين لكل \omterm{def:g7:atoms-and-molecules:molecule}{جزيء} بروبان، مقابل
$7/2 = 3.5$ \omterm{def:g8:combustion-and-fuels:complete}{للاحتراق غير التام}.

\textbf{5.} \omterm{def:g4:air-a-mixture-of-gases:air}{هواء} الخيمة المغلقة سرعان ما ينقصه الأكسجين؛ فلا يعود اللهب يحصل
على ما يكفيه منه فيحترق \omterm{def:g4:what-a-fire-needs:burning}{احتراقًا} غير تام، أصفر.

\textbf{6.} أول أكسيد الكربون: لا لون له، ولا رائحة (ولا طعم).

\textbf{7.} GHS02، اللهب: شديد القابلية للاشتعال. GHS04، قارورة
الغاز: غاز محفوظ تحت الضغط.

\textbf{8.} مثلًا: لا تستعمله إلا في \omterm{def:g4:air-a-mixture-of-gases:air}{الهواء} الطلق، ولا تستعمله أبدًا في خيمة
أو في غرفة مغلقة؛ وأبعده عن كل ما يمكن أن يحترق؛ وغيّر العبوة بعيدًا عن أي
لهب.

\textbf{9.} $132 \div 44 = 3$ مرات.

\textbf{10.} من أكسجين \omterm{def:g4:air-a-mixture-of-gases:air}{الهواء}، الذي ينتهي في ثنائي أكسيد الكربون (وفي
الماء).

\textbf{11.} \omterm{def:g8:combustion-and-fuels:fossil}{وقود أحفوري} (فهو يأتي من الغاز الطبيعي أو النفط): وحرقه يضيف
ثنائي أكسيد الكربون إلى \omterm{def:g4:air-a-mixture-of-gases:air}{الهواء}.

\textbf{12.} $450 \times 3 = \qty{1350}{g} = \qty{1.35}{kg}$ من ثنائي
أكسيد الكربون.
\end{solution}
